\documentclass[11pt]{article}
\usepackage[utf8]{inputenc}
\usepackage[margin=1in]{geometry}
\usepackage{amsmath}
\usepackage{graphicx}
\usepackage{booktabs}
\usepackage{hyperref}
\usepackage[round]{natbib}
\usepackage{float}

\title{A Translational MRI Approach to Validate Acute Axonal Damage Detection Using AxCaliber Modeling}
\author{Anonymous Authors}
\date{}

\begin{document}
\maketitle

\begin{abstract}
Axonal degeneration drives irreversible disability in multiple sclerosis (MS) and other neurodegenerative diseases, yet noninvasive detection methods lack specificity. Here we validate AxCaliber-based axonal diameter estimation as a sensitive and specific biomarker for acute axonal damage using a translational approach combining a preclinical rat model with a clinical MS cohort. In the preclinical arm ($n = 9$ rats), unilateral ibotenic acid injection into the dorsal hippocampus produced selective axonal damage in the ipsilateral fimbria without affecting myelination. AxCaliber MRI revealed a significant increase in mean axonal diameter in the injected versus control hemisphere (repeated-measures ANOVA: $F_{1,7} = 19.5$, $p = 0.003$), with histological validation showing significant correlation between neurofilament fluorescence intensity and MRI-estimated axonal diameter ($r = 0.54$, $p = 0.029$). In the clinical arm (10 relapsing-remitting MS patients, 6 healthy controls), whole-brain AxCaliber mapping revealed diffuse increases in mean axonal diameter across normal-appearing white matter (NAWM) in MS patients ($p < 0.05$, corrected), with elevated normalized index of asymmetry in MS lesions. Axonal diameter showed a negative association with disease duration, consistent with early axonal swelling followed by progressive axonal loss. These findings establish AxCaliber as a validated, specific biomarker for axonal pathology and reveal widespread diffuse axonal damage in MS NAWM.
\end{abstract}

\section{Introduction}

Axonal degeneration is the primary pathological substrate of irreversible neurological disability in multiple sclerosis (MS) and other neurodegenerative diseases \citep{trapp1998axonal, dutta2011mechanisms}. While inflammatory demyelination is the hallmark of MS relapse, progressive disability accumulation is driven by axonal loss, which can occur both within and outside visible lesions \citep{lassmann2012progressive, kutzelnigg2005cortical}. Detecting axonal damage early, before irreversible loss occurs, is therefore a critical unmet need in MS management \citep{compston2008multiple, filippi2012mri}.

Conventional diffusion MRI metrics such as fractional anisotropy (FA) and mean diffusivity (MD) are sensitive to microstructural changes but cannot distinguish axonal damage from demyelination, edema, or other pathological processes \citep{alexander2010diffusion, budde2009quantification}. Advanced diffusion models including NODDI \citep{zhang2012noddi} and diffusional kurtosis imaging \citep{fieremans2011white} provide improved specificity for neurite density but do not directly measure axonal caliber.

AxCaliber is a multicompartment diffusion MRI model that separates intra-axonal from extra-axonal water compartments, enabling estimation of axon diameter distributions \citep{assaf2008axcaliber, barazany2009vivo}. Prior studies have applied AxCaliber to the corpus callosum in MS, reporting enlarged axonal diameters \citep{huang2015corpus}. However, these studies lacked independent histological validation, and the original AxCaliber formulation is restricted to single-fiber-orientation voxels, which constitute only approximately 30\% of brain white matter.

Two critical gaps remain. First, the sensitivity and specificity of AxCaliber for detecting acute axonal damage --- as opposed to other microstructural changes --- has not been established with independent histological validation in a controlled preclinical model. Second, the extent of diffuse axonal pathology in normal-appearing white matter (NAWM) of MS patients has not been mapped using whole-brain axonal diameter estimation.

Here we address both gaps using a translational approach. In the preclinical arm, we use ibotenic acid injection to produce selective axonal damage in the rat fimbria, followed by AxCaliber MRI and immunohistochemical validation. In the clinical arm, we apply a modified whole-brain AxCaliber protocol to map axonal diameter across the white matter skeleton of MS patients and healthy controls.

\section{Results}

\subsection{Preclinical Model: Ibotenic Acid Produces Selective Axonal Damage}

Unilateral stereotaxic injection of ibotenic acid (2.5~$\mu$g/$\mu$L, 1~$\mu$L) into the right dorsal hippocampus of $n = 9$ rats produced neuronal soma death in the injected hippocampus and subsequent Wallerian-like degeneration in the ipsilateral fimbria (Table~\ref{tab:injection}).

\begin{table}[H]
\centering
\caption{Stereotaxic injection parameters for the preclinical rat model.}
\label{tab:injection}
\begin{tabular}{lr}
\toprule
\textbf{Parameter} & \textbf{Value} \\
\midrule
Neurotoxin & Ibotenic acid (NMDA agonist) \\
Concentration & 2.5 $\mu$g/$\mu$L \\
Volume & 1 $\mu$L \\
Injection site & Right dorsal hippocampus \\
Coordinates (Bregma) & $-3.8$ mm AP, 2.0 mm ML, 3.0 mm DV \\
Control injection & Saline (contralateral) \\
Survival time & 14 days \\
Sample size & $n = 9$ \\
\bottomrule
\end{tabular}
\end{table}

\subsection{AxCaliber MRI Detects Axonal Diameter Increase}

AxCaliber MRI performed 14 days post-injection revealed a significant increase in mean axonal diameter in the injected versus control fimbria (Table~\ref{tab:mri_results}).

\begin{table}[H]
\centering
\caption{Repeated-measures ANOVA results for mean axonal diameter along the fimbria tract ($n = 9$ rats). The tract was divided into 100 equidistant steps along the antero-posterior axis.}
\label{tab:mri_results}
\begin{tabular}{lrrr}
\toprule
\textbf{Factor} & \textbf{$F$-statistic} & \textbf{df} & \textbf{$p$-value} \\
\midrule
Injection (injected vs.\ control) & 19.5 & 1, 7 & \textbf{0.003} \\
Position along tract & 219.5 & 98, 686 & $<0.001$ \\
Injection $\times$ Position interaction & 11.2 & 98, 686 & \textbf{0.012} \\
\bottomrule
\end{tabular}
\end{table}

The significant injection $\times$ position interaction indicated that the axonal diameter increase was not uniform along the tract but was primarily localized posterior to the injection site, consistent with anterograde Wallerian degeneration.

\subsection{Histological Validation Confirms Axonal Specificity}

Immunohistochemistry using three markers confirmed that the ibotenic acid model produced selective axonal damage without affecting myelination (Table~\ref{tab:ihc}).

\begin{table}[H]
\centering
\caption{Immunohistochemistry results comparing injected versus control hemispheres. All comparisons use paired $t$-tests ($n = 9$).}
\label{tab:ihc}
\begin{tabular}{llllr}
\toprule
\textbf{Stain} & \textbf{Target} & \textbf{Region} & \textbf{Finding} & \textbf{$p$-value} \\
\midrule
NeuN & Neuronal somas & Hippocampus & Lower in injected & \textbf{0.026} \\
Neurofilament & Axonal integrity & Fimbria & Higher in injected & \textbf{0.047} \\
MBP & Myelin basic protein & Fimbria & No difference & 0.614 \\
\bottomrule
\end{tabular}
\end{table}

The NeuN results confirmed successful neurotoxic ablation of hippocampal neurons. Neurofilament staining showed increased intensity in the injected fimbria, consistent with axonal swelling and varicosity formation. Critically, MBP staining showed no significant difference, confirming that at 14 days post-injection, myelin remained intact despite axonal damage.

\subsection{MRI--Histology Correlation}

Neurofilament fluorescence intensity significantly correlated with AxCaliber-estimated mean axonal diameter across both injected and control fimbrias (Table~\ref{tab:correlation}).

\begin{table}[H]
\centering
\caption{Correlation between histological markers and MRI-estimated axonal diameter across all fimbria samples ($n = 18$, 9 injected + 9 control).}
\label{tab:correlation}
\begin{tabular}{lrrr}
\toprule
\textbf{Comparison} & \textbf{Pearson $r$} & \textbf{$R^2$} & \textbf{$p$-value} \\
\midrule
Neurofilament vs.\ MRI axonal diameter & 0.54 & 0.29 & \textbf{0.029} \\
MBP vs.\ MRI axonal diameter & 0.11 & 0.01 & 0.673 \\
NeuN vs.\ MRI axonal diameter & $-0.38$ & 0.14 & 0.124 \\
\bottomrule
\end{tabular}
\end{table}

The significant positive correlation between neurofilament intensity and MRI-estimated axonal diameter ($r = 0.54$, $p = 0.029$) validates that AxCaliber measurements reflect underlying axonal morphological changes. The absence of correlation with MBP confirms that the MRI measure is specific to axonal rather than myelin pathology.

\subsection{Clinical MS Cohort Demographics}

Table~\ref{tab:demographics} presents the clinical characteristics of the MS and control cohorts.

\begin{table}[H]
\centering
\caption{Demographics and clinical characteristics of the MS and healthy control cohorts.}
\label{tab:demographics}
\begin{tabular}{lrrr}
\toprule
\textbf{Characteristic} & \textbf{MS Patients} & \textbf{Controls} & \textbf{$p$-value} \\
\midrule
$N$ & 10 & 6 & -- \\
Age range (years) & 27--59 & 23--53 & 0.093$^a$ \\
Sex (female/male) & 7/3 & 3/3 & -- \\
Disease subtype & RRMS & -- & -- \\
\bottomrule
\multicolumn{4}{l}{\footnotesize $^a$ Mann-Whitney $U$ test.}
\end{tabular}
\end{table}

\subsection{Whole-Brain NAWM Shows Diffuse Axonal Diameter Elevation}

Tract-based spatial statistics (TBSS) \citep{smith2006tract} comparing NAWM axonal diameter between MS patients and healthy controls revealed widespread, mostly symmetrical increases in mean axonal diameter in MS (Table~\ref{tab:tbss}).

\begin{table}[H]
\centering
\caption{White matter tracts showing significantly elevated axonal diameter in MS NAWM compared to healthy controls ($p < 0.05$, nonparametric permutation-corrected).}
\label{tab:tbss}
\begin{tabular}{lll}
\toprule
\textbf{Tract} & \textbf{Elevated in MS} & \textbf{Laterality} \\
\midrule
Corpus callosum & Yes & Bilateral \\
Internal capsule & Yes & Bilateral \\
Corona radiata & Yes & Bilateral \\
Thalamic radiation & Yes & Bilateral \\
Inferior longitudinal fasciculus & Yes & Bilateral \\
Cingulum & Yes & Bilateral \\
Fornix & Yes & Bilateral \\
Superior longitudinal fasciculus & Yes & Bilateral \\
Inferior fronto-occipital fasciculus & Yes & Bilateral \\
Uncinate fasciculus & Yes & Bilateral \\
Tapetum & Yes & Bilateral \\
\bottomrule
\end{tabular}
\end{table}

The widespread nature of axonal diameter elevation across virtually all major white matter tracts indicates diffuse axonal pathology in NAWM, extending well beyond visible lesion boundaries.

\subsection{Lesion Analysis: Normalized Index of Asymmetry}

To quantify axonal damage within MS lesions, we computed the normalized index of asymmetry (NIA) between lesion voxels and contralateral NAWM (Table~\ref{tab:nia}).

\begin{table}[H]
\centering
\caption{Normalized index of asymmetry (NIA) for mean axonal diameter comparing lesion versus contralateral NAWM in MS patients, and the same anatomical regions in healthy controls.}
\label{tab:nia}
\begin{tabular}{lrrr}
\toprule
\textbf{Group} & \textbf{Mean NIA} & \textbf{SD} & \textbf{Significance} \\
\midrule
MS patients (lesion vs.\ contralateral) & 0.142 & 0.067 & $p < 0.01$ \\
Healthy controls (same anatomical region) & 0.023 & 0.018 & Reference \\
\midrule
MS vs.\ Controls & \multicolumn{3}{c}{$p = 0.002$ (Mann-Whitney $U$)} \\
\bottomrule
\end{tabular}
\end{table}

The NIA was significantly elevated in MS patients ($p = 0.002$), confirming greater axonal damage within lesions compared to contralateral tissue and exceeding the normal anatomical asymmetry observed in healthy controls.

\subsection{Association with Disease Duration}

Axonal diameter showed a negative association with disease duration in MS patients, observed in both TBSS and ROI-based analyses (Table~\ref{tab:duration}).

\begin{table}[H]
\centering
\caption{Association between mean axonal diameter and disease duration in MS patients. Negative correlations indicate decreasing axonal diameter with longer disease duration.}
\label{tab:duration}
\begin{tabular}{lrrr}
\toprule
\textbf{Analysis} & \textbf{Pearson $r$} & \textbf{$R^2$} & \textbf{$p$-value} \\
\midrule
TBSS skeleton (global) & $-0.61$ & 0.37 & \textbf{0.032} \\
Corpus callosum ROI & $-0.58$ & 0.34 & \textbf{0.041} \\
Corona radiata ROI & $-0.54$ & 0.29 & 0.058 \\
\bottomrule
\end{tabular}
\end{table}

This negative association is consistent with a biphasic model of axonal pathology: early disease stages feature axonal swelling (varicosities and spheroids) that increases mean diameter, followed by progressive loss of swollen axons, which reduces the mean diameter as predominantly smaller, intact axons remain \citep{nikic2011protective}.

\section{Discussion}

Our translational approach provides the first histologically validated evidence that AxCaliber MRI is both sensitive and specific for detecting acute axonal damage in vivo. The preclinical model demonstrated significant axonal diameter increases following selective neurotoxic injury ($p = 0.003$), with direct correlation to neurofilament pathology ($r = 0.54$, $p = 0.029$) and no confounding myelin changes. The clinical arm revealed that this validated biomarker detects widespread diffuse axonal pathology throughout the NAWM of MS patients.

The specificity of our preclinical model is a key strength. Ibotenic acid selectively kills neuronal somas via NMDA receptor-mediated excitotoxicity without directly affecting glial cells or myelinating oligodendrocytes \citep{paxinos2007rat}. The resulting Wallerian-like degeneration in the fimbria produces axonal pathology (confirmed by neurofilament staining) while preserving myelin (confirmed by MBP staining). This dissociation is essential because conventional diffusion MRI metrics cannot distinguish axonal from myelin pathology \citep{alexander2010diffusion}.

The clinical finding of diffuse axonal diameter elevation across virtually all major white matter tracts in MS NAWM has important implications. First, it suggests that axonal damage in MS is not confined to visible lesions but extends throughout the white matter, consistent with histopathological evidence of diffuse axonal loss in MS \citep{kutzelnigg2005cortical}. Second, the ability to detect this pathology noninvasively opens possibilities for monitoring disease progression and treatment response.

The negative association between axonal diameter and disease duration provides insight into the temporal evolution of axonal pathology. In early disease, acute axonal injury produces swelling and varicosity formation, increasing the mean diameter \citep{nikic2011protective}. Over time, selectively vulnerable smaller axons are lost, and the remaining swollen axons undergo degeneration, leading to a decrease in mean diameter. This biphasic model reconciles apparently contradictory findings in the literature and suggests that AxCaliber measurements may reflect different pathological processes depending on disease stage.

Several limitations should be noted. The clinical cohort is small ($n = 10$ MS patients, $n = 6$ controls), limiting statistical power and generalizability. The cross-sectional design cannot establish causal relationships between axonal diameter changes and disability progression. Additionally, the modified whole-brain AxCaliber protocol requires further validation against histology in crossing-fiber regions. Future longitudinal studies with larger cohorts will be needed to establish AxCaliber as a clinically useful biomarker for treatment monitoring \citep{thompson2018diagnosis}.

\section{Methods}

\subsection{Animals and Surgery}

Adult rats ($n = 9$) received unilateral stereotaxic injection of ibotenic acid (2.5~$\mu$g/$\mu$L, 1~$\mu$L) into the right dorsal hippocampus at coordinates relative to bregma: anteroposterior $-3.8$~mm, mediolateral 2.0~mm from midline, dorsoventral 3.0~mm \citep{paxinos2007rat}. Contralateral hemispheres received saline injection as internal controls. Rats survived 14 days post-surgery before MRI scanning and perfusion.

\subsection{AxCaliber MRI Protocol}

AxCaliber is a multicompartment diffusion MRI model that estimates axonal diameter distributions from diffusion-weighted signals acquired at multiple $b$-values and diffusion times \citep{assaf2008axcaliber}. The model separates intra-axonal from extra-axonal water compartments, with the intra-axonal signal modeled as restricted diffusion within cylinders of varying diameter. For the preclinical arm, fimbria tractography was performed using DTI-derived fiber orientations, and AxCaliber parameters were estimated along 100 equidistant tract steps. For the clinical arm, a modified AxCaliber model incorporating fiber orientation distributions was applied to handle crossing-fiber voxels \citep{stikov2015vivo}.

\subsection{Immunohistochemistry}

Following MRI, rats were perfused transcardially and brains processed for immunohistochemistry using three markers: NeuN (neuronal nuclei, hippocampal somas), neurofilament (axonal integrity, fimbria), and myelin basic protein (MBP, myelination status, fimbria). Fluorescence intensity was quantified in matched regions of interest across hemispheres.

\subsection{Clinical Cohort}

Ten patients with relapsing-remitting MS (age 27--59 years, 7 females) and six age-matched healthy controls (age 23--53 years, 3 females) underwent whole-brain AxCaliber MRI \citep{lublin1996defining}. MS diagnosis was established according to McDonald criteria \citep{thompson2018diagnosis}. White matter lesions were segmented from FLAIR images, and lesion masks were used to define NAWM regions for TBSS analysis.

\subsection{Statistical Analysis}

Preclinical data were analyzed using repeated-measures ANOVA (injection $\times$ position along tract) with post-hoc paired $t$-tests corrected for multiple comparisons \citep{winkler2014permutation}. MRI--histology correlations used Pearson correlation. Clinical TBSS analysis used nonparametric permutation testing with threshold-free cluster enhancement \citep{smith2006tract}. Age comparison between groups used Mann-Whitney $U$ test. All tests were two-tailed with $\alpha = 0.05$.

\section{Conclusion}

This translational study establishes AxCaliber MRI as a validated, histologically confirmed biomarker for acute axonal damage. The preclinical model demonstrates sensitivity and specificity for axonal pathology, while the clinical arm reveals diffuse axonal damage throughout the normal-appearing white matter of MS patients. Together, these findings support AxCaliber as a promising tool for noninvasive monitoring of axonal pathology in neurological disease.

\bibliographystyle{plainnat}
\bibliography{references}

\end{document}
